\section{From harmonic coefficients to pixels, and back again}
\label{sec:sht-teaching}

\paragraph{What the computer has to accomplish.}
Imagine first drawing a smooth Gaussian field on the entire sky. We then
change each pixel with a nonlinear function to reproduce a non-Gaussian
one-point distribution. Finally, we need harmonic coefficients again to
convert convergence into shear. These are three different operations. An
ordinary fast Fourier transform (FFT) helps with part of the work, but it
does not perform the whole spherical transform. This section explains why,
and follows the operations in \texttt{gpu\_scalar\_sht.py} and
\texttt{gpu\_map\_generation.py}. The mathematics applies to one source bin;
the generator repeats it for four correlated bins.

\paragraph{The three spaces we must not confuse.}
A pixel array gives the field at specified positions. A harmonic array
gives amplitudes of patterns on the sphere. A ring array is an intermediate
representation: it gives longitude Fourier amplitudes separately at each
latitude. The calculation passes through the ring array on its way between
pixels and harmonics. In our example there are
\begin{equation}
 N_{\rm pix}=12N_{\rm side}^{2}=3\,145\,728
 \qquad (N_{\rm side}=512).
\end{equation}
The chosen maximum harmonic degree is
\begin{equation}
 L=1024.
\end{equation}
These choices specify a finite numerical model. Increasing the number of
maps does not increase its angular resolution.

\subsection{Longitude waves and latitude functions}

\paragraph{Coordinates and a separated basis.}
Let $\theta$ be colatitude, measured from the north pole, and $\phi$ be
longitude. A scalar spherical harmonic separates into a latitude factor
and a longitude wave:
\begin{equation}
 Y_{\ell m}(\theta,\phi)=\lambda_{\ell m}(\theta)e^{im\phi}.
\end{equation}
We use the orthonormal convention, with the Condon--Shortley sign included
in the associated Legendre function:
\begin{equation}
 \lambda_{\ell m}(\theta)=
 \sqrt{\frac{2\ell+1}{4\pi}\frac{(\ell-m)!}{(\ell+m)!}}
 P_\ell^m(\cos\theta).
\end{equation}
Here $\ell$ controls the total angular structure, and $m$ counts longitude
oscillations. For example, $m=0$ has no longitude oscillation at all.
The first three such latitude functions are
\begin{equation}
 \lambda_{00}=\frac{1}{\sqrt{4\pi}},\qquad
 \lambda_{10}=\sqrt{\frac{3}{4\pi}}\cos\theta,\qquad
 \lambda_{20}=\sqrt{\frac{5}{16\pi}}(3\cos^2\theta-1).
\end{equation}
The constant mode describes a uniform sky. The next mode changes sign
between north and south. The third distinguishes polar and equatorial
regions. Their normalization makes the integral of $|Y_{\ell m}|^2$ over
the sphere equal to one; it is not an arbitrary amplitude choice.

\paragraph{Why only half the complex coefficients are stored.}
A real field cannot have independent positive and negative longitude
frequencies. In this convention its coefficients obey
\begin{equation}
 a_{\ell,-m}=(-1)^m a_{\ell m}^{*}.
\end{equation}
Combining the negative and positive terms therefore gives the real synthesis
formula
\begin{equation}
 f(\theta,\phi)=\sum_{\ell=0}^{L}a_{\ell0}\lambda_{\ell0}(\theta)
 +2\operatorname{Re}\sum_{m=1}^{L}\sum_{\ell=m}^{L}
 a_{\ell m}\lambda_{\ell m}(\theta)e^{im\phi}.
\label{eq:real-synthesis}
\end{equation}
The factor two restores the negative-frequency contribution. It must not
multiply the $m=0$ term. This explains both the packed coefficient array
and a normalization factor that will reappear in the iterative solver.

\paragraph{First sum over latitude functions.}
At ring $r$, all pixels have the same colatitude $\theta_r$. Define its
longitude amplitude for order $m$ by
\begin{equation}
 b_{rm}=\sum_{\ell=m}^{L}a_{\ell m}\lambda_{\ell m}(\theta_r).
\end{equation}
For $L=2$, the three stored amplitudes are
\begin{equation}
 b_{r0}=a_{00}\lambda_{00}(\theta_r)
       +a_{10}\lambda_{10}(\theta_r)+a_{20}\lambda_{20}(\theta_r).
\end{equation}
The next order contains fewer terms because $\ell$ cannot be smaller than $m$:
\begin{equation}
 b_{r1}=a_{11}\lambda_{11}(\theta_r)+a_{21}\lambda_{21}(\theta_r).
\end{equation}
The last order contains just one term:
\begin{equation}
 b_{r2}=a_{22}\lambda_{22}(\theta_r).
\end{equation}
These are weighted sums, not ordinary FFTs. The weights depend on latitude,
degree and order. A GPU can parallelize them, but an FFT library alone does
not supply this stage.

\paragraph{How the latitude weights are computed without huge factorials.}
Direct factorial formulas become numerically awkward at large degree.
Instead, the implementation starts from the diagonal function and advances
through degrees using a normalized recurrence. Put $x=\cos\theta$ and define
\begin{equation}
 A_{\ell m}=\sqrt{\frac{\ell^2-m^2}{4\ell^2-1}}.
\end{equation}
For $\ell>m$, the recurrence is
\begin{equation}
 \lambda_{\ell m}
 =\frac{x\lambda_{\ell-1,m}-A_{\ell-1,m}\lambda_{\ell-2,m}}
 {A_{\ell m}}.
\end{equation}
The starting value is
\begin{equation}
 \lambda_{mm}=(-1)^m
 \sqrt{\frac{2m+1}{4\pi}}\,
 \frac{\sqrt{(2m)!}}{2^m m!}(\sin\theta)^m.
\end{equation}
The preceding value is taken as zero at the first step. In the $m=0$
case, $A_{10}=1/\sqrt{3}$, so the first update gives
\begin{equation}
 \lambda_{10}=\frac{x/\sqrt{4\pi}}{1/\sqrt{3}}
              =\sqrt{\frac{3}{4\pi}}x.
\end{equation}
At the next step $A_{20}=2/\sqrt{15}$. Substitution gives
\begin{equation}
 \lambda_{20}=\frac{\sqrt{3}x^2-1/\sqrt{3}}
 {\sqrt{4\pi}}\frac{\sqrt{15}}{2}.
\end{equation}
Factoring the numerator makes the familiar quadratic visible:
\begin{equation}
 \lambda_{20}=\frac{\sqrt{5}}{2\sqrt{4\pi}}(3x^2-1).
\end{equation}
Thus the recurrence reproduces the elementary functions already written
above, while also extending to the thousandth degree used by the program.
Near a pole the diagonal seed can be extremely small. The kernel retains
an explicit binary exponent so that this seed is not rounded to zero
before later recurrence steps can make it significant. A leading float
and a second rounding-correction float reduce accumulated recurrence
error; this is not the same as native float64 arithmetic.

\subsection{What the longitude FFT actually does}

\paragraph{One ring is an ordinary Fourier problem.}
A ring with $N_r$ pixels samples longitudes
\begin{equation}
 \phi_{rj}=\phi_{r0}+\frac{2\pi j}{N_r},\qquad j=0,\ldots,N_r-1.
\end{equation}
After the latitude sum, synthesis only has to evaluate the longitude
waves at these positions. The phase $e^{im\phi_{r0}}$ accounts for the
starting longitude. Omitting it rotates the ring relative to the sky.
For a ring with four pixels and zero starting longitude, take a constant
amplitude one and a real positive-frequency amplitude one half:
\begin{equation}
 f_j=1+2\operatorname{Re}\left[\frac12 e^{2\pi i j/4}\right]
    =1+\cos(\pi j/2).
\end{equation}
Evaluating the four cosines gives the pixel array
\begin{equation}
 (f_0,f_1,f_2,f_3)=(2,1,0,1).
\end{equation}
Under the forward FFT convention
\begin{equation}
 F_k=\sum_{j=0}^{3}f_j e^{-2\pi i jk/4},
\end{equation}
these pixels have coefficients
\begin{equation}
 (F_0,F_1,F_2,F_3)=(4,2,0,2).
\end{equation}
The inverse FFT divides by four. Consequently, if we put the wave
amplitudes $(1,1/2,0,1/2)$ into an inverse FFT, we must multiply its result
by four to recover the intended pixels. The production code applies the
corresponding factor $N_r$ for each ring.

\paragraph{Why polar rings need special care.}
HEALPix rings have different numbers of pixels. On a ring of length $N_r$,
the sampled waves with orders $m$ and $m+N_r$ are identical after accounting
for the ring's starting phase, because
\begin{equation}
 e^{2\pi i(m+N_r)j/N_r}=e^{2\pi imj/N_r}e^{2\pi ij}
                       =e^{2\pi imj/N_r}.
\end{equation}
Their contributions must therefore be added into the same FFT bin. This
is aliasing of the sampled ring, not permission to discard either harmonic
mode. For example, on a four-pixel ring, orders one and five occupy the
same sampled frequency. Negative orders must also be included, with zero
counted once. The latitude dependence still distinguishes these modes
across the full sky. Padding a short ring and treating the extra entries
as new sky pixels would change the geometry and would be wrong.

\paragraph{Using power-of-two FFTs without changing the ring.}
Some ring lengths are inconvenient for a GPU FFT. Bluestein's identity
rewrites an exact length-$N$ discrete Fourier transform as a convolution.
Start with the elementary square expansion
\begin{equation}
 (j-k)^2=j^2-2jk+k^2.
\end{equation}
Moving the terms and dividing by two gives
\begin{equation}
 jk=\frac{j^2+k^2-(j-k)^2}{2}.
\end{equation}
The Fourier phase consequently factors as
\begin{equation}
 e^{-2\pi ijk/N}
 =e^{-\pi ij^2/N}e^{-\pi ik^2/N}e^{+\pi i(j-k)^2/N}.
\end{equation}
Define the two chirp-weighted sequences
\begin{equation}
 u_j=f_j e^{-\pi ij^2/N},\qquad v_t=e^{+\pi it^2/N}.
\end{equation}
Inserting these definitions into the original sum gives
\begin{equation}
 F_k=e^{-\pi ik^2/N}\sum_{j=0}^{N-1}u_j v_{k-j}.
\end{equation}
The remaining sum is a convolution. A padded FFT of length at least
$2N-1$, with the negative-lag chirp placed correctly, computes it without
circular wraparound contaminating the required entries. This is padding
an auxiliary convolution, not adding sky pixels. Rings requiring the same
power-of-two convolution length can share a GPU batch. Their individual
lengths, phases and normalizations remain unchanged.

\subsection{Why going back is an iterative problem}

\paragraph{The nonlinear mapping creates new angular structure.}
Even if a Gaussian input contains only low frequencies, a nonlinear
function of it generally does not. Consider the simple one-dimensional
example
\begin{equation}
 g(\phi)=\cos\phi,\qquad f(\phi)=g(\phi)^2.
\end{equation}
Using the double-angle identity gives
\begin{equation}
 f(\phi)=\frac12+\frac12\cos(2\phi).
\end{equation}
The transformed field contains a new frequency two. A model retaining
only frequencies zero and one cannot reproduce it exactly. The same issue
occurs when a nonlinear point function is applied to a band-limited sky:
it generates angular power beyond the retained $L$. The inverse task is
therefore to find the best representation at the chosen resolution.

\paragraph{Write synthesis as a matrix, without constructing that matrix.}
Let $S$ denote the linear operation that synthesizes pixels from packed
harmonic coefficients. For a trial coefficient vector $a$, its predicted
map is $Sa$. Given the transformed map $f$, the target used here is
\begin{equation}
 \min_a J(a),\qquad J(a)=\frac12\sum_p[(Sa)_p-f_p]^2.
\end{equation}
This is an unweighted full-sky pixel least-squares problem. All HEALPix
pixels have equal area. The DES footprint is applied after generating the
full-sky shear; it is not inserted into this solve.

\paragraph{An adjoint reverses a scalar product, not the transform.}
The adjoint $S^\dagger$ is defined by the identity
\begin{equation}
 \langle Sa,f\rangle_{\rm pix}=\langle a,S^\dagger f\rangle_{\rm harm}.
\end{equation}
In ordinary real matrix coordinates this is a transpose. It reverses the
order of the forward operations and accumulates contributions back into
coefficients. It is an inverse only when the relevant discrete basis is
exactly orthonormal. For a very small example, take
\begin{equation}
 S=\begin{pmatrix}1\\2\end{pmatrix},\qquad
 f=\begin{pmatrix}1\\2\end{pmatrix}.
\end{equation}
The map was generated by the coefficient $a=1$. Nevertheless, applying the
transpose gives
\begin{equation}
 S^T f=1\times1+2\times2=5.
\end{equation}
The missing normalization is supplied by the normal matrix:
\begin{equation}
 S^TS=1^2+2^2=5.
\end{equation}
Solving $5a=5$ recovers one. Thus even this one-coefficient example shows
why ``apply the adjoint'' and ``recover the coefficients'' are different
instructions.

\paragraph{The packed harmonic scalar product.}
Because positive $m$ represents both frequency signs, the harmonic inner
product must include their multiplicity:
\begin{equation}
 \langle a,b\rangle_{\rm harm}
 =\sum_{\ell}a_{\ell0}b_{\ell0}
 +2\operatorname{Re}\sum_{m>0}\sum_{\ell\ge m}a_{\ell m}^{*}b_{\ell m}.
\end{equation}
The $m=0$ coefficients are real. This is the usual real-field scalar
product expressed in compressed complex storage. With this metric, the
pixel adjoint returns
\begin{equation}
 (S^\dagger f)_{\ell m}=\sum_p f_p Y_{\ell m}^{*}(\theta_p,\phi_p).
\end{equation}
The code first performs ring Fourier sums, then latitude sums. Testing the
scalar-product identity catches phase, factor-two and normalization errors
that a timing test cannot reveal.

\paragraph{Derive the normal equations.}
Change the coefficients by a small arbitrary vector $\delta a$. The change
in the predicted map is $S\delta a$. To first order the objective changes by
\begin{equation}
 \delta J=\langle Sa-f,S\delta a\rangle_{\rm pix}.
\end{equation}
Move synthesis across the scalar product using the adjoint identity:
\begin{equation}
 \delta J=\langle S^\dagger(Sa-f),\delta a\rangle_{\rm harm}.
\end{equation}
At the minimum this must vanish for every allowed change. Therefore
\begin{equation}
 S^\dagger S a=S^\dagger f.
\end{equation}
This is a large linear system, but we need not build its matrix. Applying
it to a trial vector means synthesizing a map and then applying the
adjoint. That is the origin of the repeated forward/adjoint transforms.

\paragraph{Why pixel area appears in the implementation.}
Let $\Omega_{\rm pix}=4\pi/N_{\rm pix}$. Multiplying both sides by the same
positive number leaves the solution unchanged. The code defines
\begin{equation}
 H=\Omega_{\rm pix}S^\dagger S,\qquad
 b=\Omega_{\rm pix}S^\dagger f.
\end{equation}
It then solves $Ha=b$. For a well-sampled spherical basis, $H$ is close to
the identity, since the pixel sum times pixel area approximates the
continuous orthonormality integral. It is not exactly the identity at
finite resolution. The area factor scales the system; by itself a scalar
rescaling does not improve its condition number.

\subsection{Conjugate gradients, with every update explained}

\paragraph{Residual and search direction.}
The Metal implementation uses conjugate gradients on this normal system.
This differs from ducc0's iterative algorithm while targeting the same
least-squares problem. Start from zero coefficients:
\begin{equation}
 a_0=0,\qquad r_0=b-Ha_0=b,\qquad p_0=r_0.
\end{equation}
Here $r$ is the normal-equation residual and $p$ is a proposed direction
in coefficient space. For each step, first compute $Hp_k$ using one
synthesis and one adjoint. The distance to move along the direction is
\begin{equation}
 \alpha_k=\frac{\langle r_k,r_k\rangle_{\rm harm}}
 {\langle p_k,Hp_k\rangle_{\rm harm}}.
\end{equation}
Update the coefficients by this distance:
\begin{equation}
 a_{k+1}=a_k+\alpha_k p_k.
\end{equation}
The corresponding residual update subtracts the change in the left-hand
side of the linear system:
\begin{equation}
 r_{k+1}=r_k-\alpha_k Hp_k.
\end{equation}
If this is sufficiently small, stop. Otherwise form the ratio of residual
squared norms,
\begin{equation}
 \beta_k=\frac{\langle r_{k+1},r_{k+1}\rangle_{\rm harm}}
 {\langle r_k,r_k\rangle_{\rm harm}},
\end{equation}
and use it to retain the useful part of the previous direction:
\begin{equation}
 p_{k+1}=r_{k+1}+\beta_k p_k.
\end{equation}
The directions are conjugate with respect to $H$: in exact arithmetic,
new steps do not undo the minimization achieved along earlier directions.
The required positivity is that of the normal matrix on the represented
coefficient space. Excessive harmonic resolution or an altered mask can
change that property and cannot be treated as an innocent speed setting.

\paragraph{A two-coefficient example: first iteration.}
Consider a real system small enough to compute by hand:
\begin{equation}
 H=\begin{pmatrix}2&0\\0&1\end{pmatrix},\qquad
 b=\begin{pmatrix}2\\1\end{pmatrix}.
\end{equation}
The exact answer is $(1,1)^T$. Starting from zero gives
\begin{equation}
 r_0=p_0=(2,1)^T,\qquad Hp_0=(4,1)^T.
\end{equation}
The numerator of the step length is $2^2+1^2=5$, and its denominator is
$2\times4+1\times1=9$. Hence
\begin{equation}
 \alpha_0=\frac59,\qquad a_1=(10/9,5/9)^T.
\end{equation}
Substituting into the residual update gives
\begin{equation}
 r_1=(2,1)^T-\frac59(4,1)^T=(-2/9,4/9)^T.
\end{equation}
The new squared norm is $4/81+16/81=20/81$. Dividing it by the old norm
five gives
\begin{equation}
 \beta_0=\frac{20/81}{5}=\frac4{81}.
\end{equation}
The new direction therefore becomes
\begin{equation}
 p_1=(-2/9,4/9)^T+\frac4{81}(2,1)^T=(-10/81,40/81)^T.
\end{equation}

\paragraph{Second iteration and an explicit check.}
Applying the matrix to this direction gives
\begin{equation}
 Hp_1=(-20/81,40/81)^T.
\end{equation}
The scalar product in the denominator is
\begin{equation}
 p_1^THp_1=\frac{200+1600}{6561}=\frac{200}{729}.
\end{equation}
Consequently the next step length is
\begin{equation}
 \alpha_1=\frac{20/81}{200/729}=\frac9{10}.
\end{equation}
Substitution into the coefficient update gives
\begin{equation}
 a_2=(10/9,5/9)^T+(-1/9,4/9)^T=(1,1)^T.
\end{equation}
Multiplying this answer by $H$ returns $(2,1)^T=b$, so the residual is
zero. The sky calculation follows these same scalar-product and vector
updates; the expensive operation is applying $H$ through two transforms.

\paragraph{Two residuals answer different questions.}
The solver monitors a relative normal-equation residual,
\begin{equation}
 \epsilon_{\rm normal}=
 \frac{\|b-Ha\|_{\rm harm}}{\|b\|_{\rm harm}}.
\end{equation}
A small value means we have approximately solved the chosen finite linear
system. It does not mean the reconstructed pixels equal every input pixel.
The map residual is instead
\begin{equation}
 e_{\rm map}=Sa-f.
\end{equation}
Unrepresentable high-frequency structure can remain in this map residual
even at the exact least-squares minimum. Nor does a small normal residual
alone guarantee small coefficient error for an ill-conditioned system.
We therefore compare maps and statistics against ducc0 as well. The GPU
solver recomputes $b-Ha$ independently at the end rather than trusting
only the recursively updated residual. The default target is $10^{-6}$;
the independent final check rejects values above three times that target.

\subsection{Where GPU acceleration enters, and what was tested}

\paragraph{Follow one realization through the program.}
The CPU draws the correlated Gaussian coefficients using the original
random stream and double-precision calibration. On the GPU, latitude sums
and longitude FFTs synthesize the scalar sky. The GPU applies the fitted
point transformation, then the iterative scalar inverse just derived.
The resulting convergence coefficients return to the CPU for the final
spin-two shear synthesis with ducc0. A scalar harmonic transform cannot
silently substitute for that spin-two operation.

The convergence-to-shear coefficient factor used by the existing generator
is, for $\ell\ge2$,
\begin{equation}
 q_\ell=-\sqrt{\frac{(\ell+2)(\ell-1)}{\ell(\ell+1)}}.
\end{equation}
Degrees zero and one are set to zero in this conversion. The minus sign
belongs to the implemented harmonic/shear convention; changing it without
changing the matching spin-two convention would change the maps. The
existing component signs, footprint selection and noise draws are retained.
Saved pixels are float32. Measurements decode them to float64, so their
values refer to the actual saved maps.

\paragraph{Which numerical checks support the optional path.}
At NSIDE512 and $L=1024$, the scalar synthesis and adjoint differed from
ducc0 by relative RMS values approximately $1.07\times10^{-6}$ and
$1.14\times10^{-6}$. The normalized adjoint scalar-product discrepancy was
$2.09\times10^{-11}$, and repeated outputs agreed bitwise in that test.
Two full four-source-bin noiseless realizations had maximum per-bin map
RMS differences below $2.27\times10^{-6}$. Their tested aperture moments
through sixth order differed by less than $2.87\times10^{-6}$ when scaled
by absolute moments. These checks are recorded under
\texttt{tests/gptg\_pilot/gpu\_sht/}. They test numerical agreement for the
specified resolution and calibration, not the physical adequacy of the
GPTG fit or the full TreeCorr covariance.

\paragraph{Why the scheduling uses both CPU and GPU.}
The current mixed MPI allocation has one coordinator, three CPU map
workers and one GPU scalar-transform worker. Each worker may use three
native CPU transform threads, including the GPU worker during its CPU
shear stage. Thus there are at most twelve such CPU threads; each BLAS
library remains limited to one thread. Maps are independent jobs, and
the coordinator sends the next map to a worker when that worker finishes.
The GPU's internal execution is managed by Metal, not by choosing twelve
GPU workers. Each completed map records its arithmetic backend.

A controlled 32-map comparison took 47.15 seconds with the previous CPU
float64 allocation and 38.18 seconds with the mixed allocation. Reversing
the order gave 48.12 and 38.56 seconds. These are about a 20 percent
reduction in elapsed time, or a 25 percent increase in throughput. They
include more than the scalar transform alone. The full resumed run can
have different throughput because disk activity and other jobs vary.
Previously saved maps are retained; a numerical-mode history records the
transition. Full TreeCorr measurements remain a separate HPC task.
